\section{The original canonical differential before trace normalization}

\begin{lemma}[The same native map on every original raw form]
\label{lem:kn-global-trace-native-branch-differential}
\lean{EnvelopeProperty.Vanishing.NadelGlobalCanonicalTrace.nativeCanonicalPullbackMap,EnvelopeProperty.Vanishing.NadelGlobalCanonicalTrace.nativeCanonicalPullbackMap_apply,EnvelopeProperty.Vanishing.NadelGlobalCanonicalTrace.nativeCanonicalPullbackMap_raw}
\leanok
Let $f:\operatorname{Spec}B\to\operatorname{Spec}k$ be smooth of pure
relative dimension $n$, and let $\varphi:B\to C$ be any original ring map.
Put $g=\operatorname{Spec}(\varphi)$.  There is a produced $B$-linear map
\[
 d_g:\Gamma(\omega_f)\longrightarrow
      \operatorname{Res}_{\varphi}\Gamma(\omega_{g\circ f})
\]
whose value on every original section is the original adjoint canonical
differential.  For each original raw presheaf top form $z$ on the whole
source, with original sheafification unit $\eta$, it satisfies
\[
 d_g(\eta z)=\eta(d_g^{\mathrm{raw}}z).
\]
The target object, original ring actions, raw differential, and units in
this formula are literal.  No etaleness or prescribed differential
comparison is an input.
\uses{lem:kn-global-trace-native-spectrum}
\end{lemma}
\begin{proof}
Smoothness of the source makes $\omega_f$ quasicoherent.  Start with the
actual scalar-extension adjunction unit on its native $B$-section module.
The published original-ring affine pullback tensor comparison identifies
the resulting $C$-tensor module with the sections of $g^*\omega_f$.
Apply the section functor to the actual canonical differential
$g^*\omega_f\to\omega_{g\circ f}$ and restrict through the original
$\varphi$ action.  This defines the displayed map for an arbitrary
original $g$; no differential isomorphism is asserted.

On the original unit tensor, the published affine comparison is the
same original pullback-pushforward unit.  The original adjunction triangle
therefore identifies the value with the actual adjoint canonical
differential on every source section.  The existing whole raw-unit
identity for the canonical sheaf morphism then gives the displayed raw
formula on every original $z$.  This is precisely the differential to
be used in the mixed-counit normalization, including along its branch
divisor; no substitute comparison map is introduced.
\end{proof}


\begin{lemma}[Original ordinary coordinate differentials in the native unit]
\label{lem:kn-global-trace-native-ordinary-unit}
\lean{EnvelopeProperty.Vanishing.NadelGlobalCanonicalTrace.affineCotangentModuleIso_inv_D,EnvelopeProperty.Vanishing.NadelGlobalCanonicalTrace.affineOrdinaryCotangentToOpen_D}
\leanok
Let $A$ be an original algebra over the field $k$.  Under the actual
categorical-to-ordinary coordinate cotangent comparison, the original
$D_{A/k}(a)$ maps to the categorical differential of the same original
$a$.  Consequently the actual ordinary coordinate map to any original
open of $\operatorname{Spec}A$ sends it to the native presheaf differential
of the same original coefficient section.
\end{lemma}
\begin{proof}
The categorical coordinate module uses the algebra structure obtained
from the original ring homomorphism $k\to A$.  The ordinary module uses
the given original $k$-algebra structure.  The proved algebra-structure
equality identifies these two, and this exact equality defines their
coordinate module isomorphism.  The universal differential is a dependent
function of the algebra structure.  Transport along an equality carries
its value to the value for the transported structure: eliminate that
equality, after which the statement is reflexivity.  Apply this to the
original coordinate structure equality and its inverse.  The existing
native categorical coordinate map sends every categorical $da$ to the
native presheaf differential of its same original coefficient image.
Composing gives the asserted ordinary-coordinate formula, with no
replacement algebra map or desired-unit hypothesis.
\end{proof}
